\section{Fundamentos de la energía nuclear: fisión, fusión y control}

Esta sección empieza por asimilar los conceptos básicos. La energía nuclear tiene un origen común, el defecto de masa: la masa total antes y después de la reacción no es la misma, y la masa faltante se convierte en energía según la ecuación masa-energía de Einstein. La fisión es la división de un núcleo pesado en núcleos más ligeros; la fusión es la unión de núcleos ligeros en un núcleo más pesado. Ambas pueden liberar energía, pero difieren por completo en combustible, condiciones de reacción, residuos y características de seguridad.

¿Por qué la fusión nuclear tiene que llevar el calificativo «controlada»? Porque la liberación no controlada es un proceso explosivo del tipo de la bomba de hidrógeno; lo que busca la fusión nuclear controlada es entregar energía de forma estable, con la potencia, la duración y los márgenes de seguridad que requiere el ser humano. Esa es la enorme diferencia entre «liberar energía» y «un producto energético»: lo primero demuestra que la reacción puede ocurrir; lo segundo exige que la reacción pueda sostenerse, regularse, ceder su calor, conectarse a la red y recibir mantenimiento.

\[
E=\Delta m c^2
\]

Aquí, $E$ es la energía liberada, $\Delta m$ es la masa que se pierde entre el inicio y el final de la reacción, y $c$ es la velocidad de la luz. Como $c^2$ es enorme, incluso una diferencia de masa muy pequeña corresponde a una cantidad inmensa de energía. Esta fórmula explica la densidad energética de la energía nuclear y también por qué la fusión nuclear, una vez controlada, recibe el nombre de santo grial de la energía.

\begin{knowledgebox}{Términos clave: fisión, fusión, fusión controlada}
\begin{tabularx}{\textwidth}{p{0.18\textwidth}p{0.34\textwidth}X}
\toprule
Concepto & Mecanismo central & Implicación de ingeniería \\
\midrule
Fisión & Un núcleo pesado se divide en núcleos más ligeros y libera la energía correspondiente a la diferencia de masa & Ya está comercializada, pero implica combustible radiactivo, residuos y presión regulatoria en materia de seguridad. \\
Fusión & Núcleos ligeros chocan y se funden en un núcleo más pesado, liberando energía & El potencial del combustible es enorme, pero requiere temperaturas altísimas y condiciones de confinamiento. \\
Fusión no controlada & La energía se libera en un tiempo extremadamente corto & Lógica de un arma, no de la generación eléctrica. \\
Fusión nuclear controlada & Libera energía de forma estable con la potencia de diseño & Requiere que confinamiento, control, extracción de calor, materiales, operación y mantenimiento, y el sistema de red eléctrica superen juntos el umbral. \\
\bottomrule
\end{tabularx}
\end{knowledgebox}

\begin{figure}[H]
\centering
\includegraphics[width=0.96\textwidth]{fusion-energy-chain.png}
\caption{Cadena energética de la fusión nuclear controlada: la ruta central desde el combustible y el plasma hasta la entrega a la red eléctrica. Diagrama conceptual propio, elaborado a partir de la conversación de 00:04:22--00:08:34.}
\end{figure}

\begin{knowledgebox}{Lectura de la figura: una central de fusión no termina al «encender» el plasma}
En la figura, el tramo de Fuel a Plasma es la entrada de la reacción física; Confinement consiste en evitar que el plasma a alta temperatura toque las paredes sólidas; Heat y Grid representan la conversión en producto propia de la ingeniería: la energía debe extraerse, convertirse en electricidad e integrarse de forma estable a la red. Muchas discusiones se quedan en los dos primeros pasos, pero una central comercial tiene que recorrer la cadena completa.
\end{knowledgebox}

\subsection{Resumen de la sección}

Tanto la fisión como la fusión provienen del defecto de masa, pero el reto de la fusión nuclear controlada consiste en «liberar y aprovechar la energía de forma estable, controlada y a bajo costo». No es una explosión única, sino un sistema energético que opera a largo plazo.
